\scanpage{38}
\mspage{-9. 7-}
n° 215

opère à gauche. Désignons par $\Gamma_F$ le groupe de trivialité de $F$. Supposons que $F$ et $G/\Gamma_F$ soient munis de structures d'espèce $\unsure{\mathbb E|E}$ satisfaisant à la condition suivante :

C. \uline{si} $a\rightarrow s_a^*$ \uline{est une application dans} $G/\Gamma_{\handwritten{\unsure{F2|F1}}}$ \uline{d'un ensemble} $A$ \uline{muni d'une structure d'espèce} $\mathbb E$, \uline{une condition nécessaire et suffisante pour que cette application soit une} $\mathbb E$-\uline{application est que la famille d'applications} $(\xi\rightarrow s_a^*.\xi)_{a\in A}$ \uline{de} $F$ \uline{dans lui-même soit une} $\mathbb E$-\uline{famille}.

Il résulte en particulier de là que, pour tout $s^*\in G/\Gamma_F$, l'application $\xi\rightarrow s^*.\xi$ est un automorphisme de la structure d'espèce $\mathbb E$ de $F$. Nous allons montrer que la condition C. entraîne que l'espèce de structures $\unsure{\mathbb F|\mathbb E}(F;G)$ est plus puissante que $\mathbb E$. Soit en effet $F'$ un ensemble muni d'une structure d'espèce $\mathbb F(F;G)$, et soit $\Theta$ son ensemble de repères. Si $\theta\in\Theta$ et $s^*\in G/\Gamma_F$, nous désignerons par $\theta s^*$ le repère $\xi\rightarrow\theta(s^*.\xi)$. Tout élément $\theta$ de $\Theta$ permet de transporter à $F'$ la structure d'espèce $\mathbb E$ de $F$ ; tout autre élément de $\Theta$ étant de la forme $\theta s^*$, avec \handwritten{\inserted{un}} $s^*\in G/\Gamma_F$, la structure d'espèce $\mathbb E$ ainsi définie sur $F'$ ne dépend pas du choix de $\theta$. Les morphismes de l'espèce de structures $\mathbb F(F;G)$ étant des isomorphismes, il est clair que ce sont aussi des isomorphismes pour les structures d'espèce $\mathbb E$ de ces ensembles. On pourra donc parler d'espaces fibrés d'espèce $\mathbb F(F;G)$.

Soit maintenant $H$ un autre ensemble sur lequel $G$ opère à gauche et dont le groupe de trivialité $\Gamma_H$ contienne $\Gamma_F$ ; il y a alors une application canonique $s^*\rightarrow s^{**}$ de $G/\Gamma_F$ sur $G/\Gamma_H$. Supposons que $H$ satisfasse à la condition C. Ci-dessus (où on remplace $F$ par $H$), et que de plus la condition suivante soit satisfaite :

C'. \uline{l'application canonique de} $G/\Gamma_F$ \uline{dans} $G/\Gamma_H$ \uline{est une}

% archive sheet number: 37

